\section{Change of basis, power sums, and sector bookkeeping}
\label{sec:app_bnt_power_sector}

This section collects the finite-dimensional linear algebra behind alternative
BNT comparisons, the extrapolation identities for finite sums of geometric
sequences, and the bookkeeping identities for sector decompositions.  A sector
has copy weights $\mu_{j,q}\ne0$, coefficient
$c_{N,j}=\sum_q\mu_{j,q}^N$, representative bond dimension $D_j$, and total
bond dimension $D_{\mathrm{tot}}=\sum_j r_jD_j$.

\begin{lemma}[Invertible change of basis]
    \label{lem:bnt_invertible_change_basis}
    \lean{exists_invertible_changeBasis}
    \leanok
    Let $v_1, \ldots, v_g$ and $w_1, \ldots, w_g$ be two linearly independent
    families in a complex vector space with the same span. Then there exists an
    invertible matrix $U \in \GL_g(\C)$ such that
    $w_j = \sum_{i=1}^g U_{ij} v_i$ for every $j$.
\end{lemma}

\begin{proof}\leanok
    Because the two families span the same subspace, each $w_j$ has a unique
    expansion in the basis $\{v_i\}_{i=1}^g$; these coefficients form the
    columns of $U$. If $U a = 0$, then the corresponding linear relation among
    the $w_j$ is trivial because the $w_j$ are linearly independent. Thus $U$
    is injective, hence invertible.
\end{proof}

\begin{lemma}[Eventual change of basis from equal spans]
    \label{lem:bnt_eventual_change_basis}
    \lean{MPSTensor.eventually_exists_invertible_changeBasis}
    \leanok
    \uses{lem:bnt_invertible_change_basis, def:mpv_overlap}
    If two finite families of blocks have asymptotically orthonormal overlaps
    within each family and their MPV spans agree for every system size $N$,
    then for all sufficiently large $N$ there exists an invertible matrix
    $U_N$ expressing the MPV states of one family as linear combinations of
    the MPV states of the other.
\end{lemma}

\begin{proof}\leanok\uses{lem:bnt_overlap_orthonormal_li, lem:bnt_invertible_change_basis}
    By Lemma~\ref{lem:bnt_overlap_orthonormal_li}, both families are
    linearly independent for all sufficiently large $N$. At such an $N$, the
    span equality lets us apply Lemma~\ref{lem:bnt_invertible_change_basis}.
\end{proof}

\begin{theorem}[Permutation rigidity for equal block counts]
    \label{thm:bnt_perm_same_card}
    \lean{MPSTensor.exists_perm_dimEq_gaugePhaseEquiv_of_overlapOrtho}
    \leanok
    \uses{def:injective, def:mpv_overlap, def:gauge_phase_equiv}
    Let $(A_j)_{j=1}^g$ and $(B_j)_{j=1}^g$ be injective families of positive
    bond dimension satisfying trace-preserving normalization, with self-overlaps tending to~$1$ and
    off-diagonal overlaps tending to~$0$ within each family. If the spans of
    the MPV states agree for every system size $N$, then there is a permutation
    $\pi$ of $\{1, \ldots, g\}$ such that $A_j$ and $B_{\pi(j)}$ have the
    same bond dimension and are gauge-phase equivalent.
\end{theorem}

\begin{proof}\leanok\uses{lem:bnt_eventual_change_basis, thm:overlap_decay, thm:overlap_decay_rect}
    For large $N$, Lemma~\ref{lem:bnt_eventual_change_basis} gives an
    invertible matrix $U_N$ relating the two MPV families. Fix $j$ and write
    $\ket{V^{(N)}(B_j)}=\sum_i u_{ij}^{(N)}\ket{V^{(N)}(A_i)}$. With
    $G_A(N)=(O_{A_iA_\ell}(N))_{i,\ell}$ and
    $b_j(N)=(O_{A_iB_j}(N))_i$, this gives
    \begin{align}
        G_A(N)\overline{u_j(N)} &= b_j(N), \notag\\
        O_{B_jB_j}(N) &= u_j(N)^{\top}G_A(N)\overline{u_j(N)}. \notag
    \end{align}
    If all mixed overlaps $O_{A_iB_j}(N)$ tended to $0$, then
    $b_j(N)\to0$; since $G_A(N)\to\Id$, one obtains
    $\overline{u_j(N)}\to0$ and hence $u_j(N)\to0$. Thus
    $O_{B_jB_j}(N)\to0$, contradicting $O_{B_jB_j}(N)\to1$.
    Thus each $B_j$ has a matched block $A_{f(j)}$ with non-decaying mixed
    overlap. Theorems~\ref{thm:overlap_decay_rect}
    and~\ref{thm:overlap_decay} then force equal bond dimension and
    gauge-phase equivalence for each match. Off-diagonal decay within the
    $B$-family makes $f$ injective, hence bijective, and its inverse is the
    required permutation.
\end{proof}

\begin{lemma}[Permutation rigidity with asymptotically orthonormal overlaps]\label{lem:bnt_perm_simple}
    \lean{MPSTensor.exists_eq_numBlocks_and_equiv_gaugePhase_of_overlapOrtho}
    \leanok
    \uses{def:injective, def:mpv, def:mpv_overlap, def:gauge_phase_equiv}
    Let $(A_j)_{j=1}^{g_A}$ and $(B_k)_{k=1}^{g_B}$ be two
    families of injective tensors of positive bond dimension satisfying the TP normalization
    $\sum_i (A_j^i)^\dagger A_j^i = \Id$ and
    $\sum_i (B_k^i)^\dagger B_k^i = \Id$, whose self-overlaps
    converge to~$1$ and whose cross-overlaps within each
    family converge to~$0$.
    If for every system size~$N$,
    \begin{align}
        \spn_{\C}\{\ket{V^{(N)}(A_j)} : 1 \le j \le g_A\}
        &=
        \spn_{\C}\{\ket{V^{(N)}(B_k)} : 1 \le k \le g_B\}, \notag
    \end{align}
    then $g_A = g_B$, and there is a permutation $\pi$ of the common
    block index set such that $A_j$ and $B_{\pi(j)}$ have the same bond
    dimension and are gauge-phase equivalent for each~$j$.
\end{lemma}

\begin{proof}\leanok\uses{lem:bnt_overlap_orthonormal_li, thm:bnt_perm_same_card}
    By Lemma~\ref{lem:bnt_overlap_orthonormal_li}, both families are
    linearly independent for all sufficiently large $N$. At such an $N$, the
    common span has the same dimension when computed from the $A$-family or
    the $B$-family, so $g_A = g_B$. After identifying the two index sets,
    Theorem~\ref{thm:bnt_perm_same_card} yields the permutation conclusion.
\end{proof}

\begin{lemma}[Eventual vanishing of a finite geometric sum]
    \label{lem:finite_geometric_sum_eventual_vanishing}
    \lean{MPSTensor.SectorWeightData.geom_sum_eventually_zero}
    \leanok
    Let $w_1,\ldots,w_n$ be nonzero complex numbers. If
    $\sum_i c_i w_i^N=0$ for every sufficiently large $N$, then the same
    equality holds for every non-negative exponent.
\end{lemma}

\begin{proof}\leanok
    Induct on the number of terms. Subtracting $w_1$ times the equality at
    exponent $N$ from the equality at exponent $N+1$ gives
    $\sum_{i=2}^{n} c_i(w_i-w_1)w_i^N=0$.
    The induction hypothesis yields this equality at every exponent, so the
    original sum $S_N=\sum_i c_iw_i^N$ satisfies $S_{N+1}=w_1S_N$ for every
    $N$. Since $S_N$ eventually vanishes and $w_1\ne0$, it follows that
    $S_0=0$ and hence $S_N=0$ for every $N$.
\end{proof}

\begin{lemma}[Eventual power-sum equality holds at every exponent]
    \label{lem:eventual_power_sum_equality_all_exponents}
    \lean{MPSTensor.SectorWeightData.power_sums_eq_of_eventually_eq_hetero}
    \leanok
    \uses{lem:finite_geometric_sum_eventual_vanishing}
    Let $a_1,\ldots,a_m$ and $b_1,\ldots,b_n$ be nonzero complex numbers. If
    their power sums agree for every sufficiently large exponent, then
    $\sum_i a_i^k=\sum_j b_j^k$ for every non-negative exponent $k$.
\end{lemma}

\begin{proof}\leanok
    Set $w_l=a_l$ for $l\leq m$ and $w_{m+l}=b_l$ for $l\leq n$, with
    coefficients $c_l=1$ for $l\leq m$ and $c_l=-1$ for $l>m$. Then
    $\sum_{l=1}^{m+n}c_lw_l^k=\sum_i a_i^k-\sum_j b_j^k$.
    The hypothesis makes this expression zero for every sufficiently large
    $k$. Lemma~\ref{lem:finite_geometric_sum_eventual_vanishing} makes it zero
    for every $k$, which is the required power-sum equality.
\end{proof}

\begin{theorem}[Equal power sums imply equal multisets]\label{thm:power_sum_multiset}
    \lean{Matrix.sum_pow_eq_implies_multiset_eq}
    \leanok
    Let $\alpha, \beta : \{0, \ldots, n-1\} \to \C$.
    If $\sum_i \alpha_i^k = \sum_i \beta_i^k$ for all $k \ge 1$,
    then the multisets $\{\alpha_i\}$ and $\{\beta_i\}$ are equal.
\end{theorem}

\begin{proof}
    \leanok
    \uses{thm:newton_girard}
    Construct diagonal matrices $\diag(\alpha)$ and
    $\diag(\beta)$. The power-sum hypothesis gives
    $\tr(\diag(\alpha)^k) = \tr(\diag(\beta)^k)$.
    By Theorem~\ref{thm:newton_girard}, the characteristic
    polynomials agree: $\prod_i (X - \alpha_i) = \prod_i (X - \beta_i)$.
    Extracting roots gives multiset equality.
\end{proof}

\begin{theorem}[Representative dimensions are bounded by the total bond dimension]
    \label{thm:sector_basis_dimension_le_total_dimension}
    \lean{MPSTensor.SectorDecomposition.basisDim_le_totalDim}
    \lean{MPSTensor.SectorDecomposition.basisCount_le_totalDim}
    \leanok
    \uses{def:sector_weight_data}
    Let the representative $A_j$ have bond dimension $D_j$, and let $D$ be
    the bond dimension of the flattened sector tensor. Then $D_j\leq D$.
    If every $D_j$ is positive and there are $g$ representatives, then also
    $g\leq D$.
\end{theorem}

\begin{proof}
    \leanok
    \uses{def:sector_weight_data}
    Since $r_j>0$, choose one copy of $A_j$. Its contribution $D_j$ is one
    non-negative summand in the total dimension of the flattened copy family,
    so $D_j\leq D$.

    For the second assertion, every representative has multiplicity
    $r_j\geq1$, and each copy has positive dimension $D_j>0$. Therefore
    \begin{align}
        g=\sum_j1
        \leq\sum_j r_j
        \leq\sum_j r_jD_j
        =D.
        \notag
    \end{align}
\end{proof}

\begin{theorem}[Blocked sector coefficients and assembled tensor]%
    \label{thm:sector_decomposition_block_tensor_mpv}
    \lean{MPSTensor.SectorDecomposition.coeff_blockTensor}
    \lean{MPSTensor.SectorDecomposition.sameMPV₂_blockTensor_toTensor}
    \leanok
    \uses{def:sector_decomposition_block_tensor, thm:block_tensor_to_tensor_from_blocks}
    The coefficient of the blocked decomposition satisfies
    \begin{align}
        c_{P^{[p]},j}^{(N)}
        =\sum_q(\mu_{j,q}^{p})^N
        =c_{P,j}^{(Np)}.
        \label{eq:bnt_sector_block_coeff}
    \end{align}
    Moreover, blocking the tensor assembled from $P$ gives the same MPV family
    as the tensor assembled from the blocked sector decomposition.
\end{theorem}

\begin{proof}\leanok
    \uses{thm:block_tensor_to_tensor_from_blocks}
    The coefficient identity in
    (\ref{eq:bnt_sector_block_coeff}) is $(\mu^p)^N=\mu^{pN}$. The tensor
    identity is Theorem~\ref{thm:block_tensor_to_tensor_from_blocks} applied to
    the flattened copy family.
\end{proof}

\begin{lemma}[Existence of an MPV phase class]
    \label{lem:mpv_phase_class_nonempty}
    \lean{MPSTensor.MPVPhaseClassData.g_pos_of_r_pos}
    \leanok
    \uses{def:mpv_phase_class_data}
    A nonempty finite family of blocks has at least one MPV phase class.
\end{lemma}

\begin{proof}\leanok
    Choose one original block index. Its occurrence in the phase-class
    enumeration supplies a phase class containing it.
\end{proof}

\begin{lemma}[Regrouping a matrix sum by MPV phase classes]
    \label{lem:mpv_phase_class_regroup_matrix}
    \lean{MPSTensor.MPVPhaseClassData.regroup_matrix}
    \leanok
    \uses{def:mpv_phase_class_data}
    Let $F_k$ be a finite family of square matrices indexed by the original
    blocks. If $k=(j,q)$ enumerates the copies in the MPV phase classes, then
    $\sum_j\sum_q F_{j,q}=\sum_k F_k$.
\end{lemma}

\begin{proof}\leanok
    Evaluate both sides at a pair of matrix indices and apply the bijective
    enumeration of the finite phase-class quotient.
\end{proof}

\begin{definition}[Single-copy sector decomposition]
    \label{def:trivial_sector_decomp}
    \lean{MPSTensor.trivialSectorDecomp}
    \leanok
    \uses{def:sector_weight_data}
    Given nonzero weights $(\mu_k)_{k=1}^r$ and blocks $(B_k)_{k=1}^r$, the
    \emph{single-copy sector decomposition} has $r$ sectors, one basis block
    $B_k$ per sector, multiplicity~$1$, and sector weight $\mu_k$. As a finite
    sum,
    \begin{align}
        \sum_{k=1}^{r}\mu_k^{\,N}V^{(N)}(B_k)_\sigma
         & =
        V^{(N)}\!\left(\bigoplus_k\mu_kB_k\right)_\sigma.
        \label{eq:canonical_single_copy_sector}
    \end{align}
\end{definition}

\begin{lemma}[Single-copy sector decomposition preserves the represented family]
    \label{lem:trivial_sector_decomp_same_mpv}
    \lean{MPSTensor.sameMPV₂_trivialSectorDecomp}
    \leanok
    \uses{def:trivial_sector_decomp}
    The single-copy sector decomposition of $(\mu_k,B_k)_k$ represents the same
    MPS family as $\bigoplus_k\mu_kB_k$.
\end{lemma}

\begin{proof}\leanok
    \uses{def:trivial_sector_decomp}
    Expanding the sector decomposition gives the left-hand side of
    (\ref{eq:canonical_single_copy_sector}), which is the finite block-sum
    expression for $\bigoplus_k\mu_kB_k$.
\end{proof}

\begin{lemma}[Sector-decomposition coefficient is not eventually zero]
    \label{lem:sector_decomposition_coeff_not_eventually_zero}
    \lean{MPSTensor.SectorDecomposition.coeff_not_eventually_zero}
    \leanok
    \uses{def:sector_weight_data,
        lem:finite_geometric_sum_eventual_vanishing}
    For any sector decomposition $P$ and any sector $j$, the coefficient
    $c_{P,j}^{(N)}=\sum_q\mu_{j,q}^N$ is not eventually zero in $N$.
\end{lemma}

\begin{proof}\leanok
    If the coefficient vanished for every $N\geq M$, the finite geometric-sum
    extrapolation would make it vanish at every non-negative exponent. However,
    at exponent zero, $c_{P,j}^{(0)}=\sum_q\mu_{j,q}^0=r_j>0$, which is
    impossible.
\end{proof}

\begin{lemma}[Sector coefficient is not eventually zero]
    \label{lem:sector_bnt_coeff_not_eventually_zero}
    \lean{MPSTensor.IsBNTCanonicalForm.coeff_not_eventually_zero}
    \leanok
    \uses{def:sector_bnt_cf,
        lem:sector_decomposition_coeff_not_eventually_zero}
    For a BNT canonical form $P$ and any sector $j$, the sector coefficient
    $c_{P,j}^{(N)} = \sum_q \mu_{j,q}^N$ is not eventually zero in $N$. Only
    nonvanishing of the copy weights is used, not any modulus condition.
\end{lemma}

\begin{proof}\leanok
    Apply Lemma~\ref{lem:sector_decomposition_coeff_not_eventually_zero}; its
    conclusion does not require the canonical-form assumptions.
\end{proof}

\subsection{Proportional matching and finite bijection extraction}
\label{subsec:app_bnt_proportional_matching}

The equal-MPV matching theorems in Chapter~\ref{ch:bnt} are specializations of
an eventual proportionality argument. We record that general mechanism and the
finite-set construction used to package matches in both directions.

\begin{theorem}[Exact sector matching under eventual proportionality]
    \label{lem:sector_bnt_proportional_exact_match}
    \lean{MPSTensor.exists_block_match_exact_of_eventuallyProportional}
    \leanok
    \uses{def:sector_bnt_cf, lem:sector_bnt_combined_family_eventually_li,
        lem:sector_bnt_coeff_not_eventually_zero}
    Let $P$ and $Q$ be BNT canonical forms whose total MPV families are
    eventually nonzero and proportional. For every sector $j$ of $P$, there is
    a sector $k$ of $Q$ with equal bond dimension, gauge-phase-equivalent basis
    tensor after the dimension identification, and non-decaying cross-overlap.
\end{theorem}

\begin{proof}\leanok
    \uses{lem:sector_bnt_combined_family_eventually_li,
        lem:sector_bnt_coeff_not_eventually_zero,
        thm:nt_overlap_norm_one_gauge}
    If the chosen sector decayed against every sector of $Q$, separate the
    non-decaying matched sectors from the decaying residue. The combined-family
    Gram criterion makes that residue together with the $Q$ family eventually
    linearly independent. Exact coefficient extraction would then make the
    chosen sector coefficient eventually zero, contradicting its power-sum
    nonvanishing. The resulting non-decaying overlap gives dimension equality
    and gauge-phase equivalence by normal-block overlap rigidity.
\end{proof}

\begin{theorem}[Full-basis matching under eventual proportionality]
    \label{lem:sector_bnt_proportional_full_basis_match}
    \lean{MPSTensor.forall_k_exists_j_nondecaying_overlap_of_eventuallyProportional}
    \leanok
    \uses{lem:sector_bnt_proportional_exact_match}
    Under the same hypotheses, every sector of $Q$ has an equal-dimensional,
    gauge-phase-equivalent partner in $P$ with non-decaying cross-overlap.
\end{theorem}

\begin{proof}\leanok
    Reverse the eventual proportionality and apply
    Theorem~\ref{lem:sector_bnt_proportional_exact_match} to each sector of
    $Q$; symmetry of the overlap converts the resulting witness to the stated
    orientation.
\end{proof}

\begin{lemma}[Bijection from matches in both directions]
    \label{lem:sector_bnt_bijection_from_matches}
    \lean{MPSTensor.bijection_from_matches}
    \leanok
    \uses{def:sector_bnt_cf}
    Suppose every sector of $Q$ has an equal-dimensional,
    gauge-phase-equivalent partner in $P$ with non-decaying overlap, and the
    converse holds from $P$ to $Q$. Then the forward partners may be chosen as
    a bijection between the two finite sector sets, preserving all three
    properties.
\end{lemma}

\begin{proof}\leanok
    Basis separation makes each partner map injective: two sectors with the
    same partner would be gauge-phase equivalent to one another. The injections
    in both directions force equal finite cardinalities, so the forward map is
    bijective.
\end{proof}
